\section{Architecture}

\label{sec:architecture}

\subsection{Lux C++ EVM}

The Lux EVM is a ground-up C++ reimplementation of the Ethereum Virtual Machine, designed for GPU acceleration~\cite{lux-evmgpu-benchmark}. Key architectural differences from geth/evmone:

\begin{itemize}
    \item \textbf{Opcode compilation}: EVM bytecode is compiled to an intermediate representation (IR) suitable for GPU dispatch, rather than interpreted instruction-by-instruction.
    \item \textbf{Parallel state access}: A transactional memory model allows concurrent reads with conflict detection, enabling multiple transactions to execute simultaneously when they access disjoint state.
    \item \textbf{GPU shaders}: Arithmetic opcodes (\texttt{ADD}, \texttt{MUL}, \texttt{EXP}, \texttt{KECCAK256}) are implemented as compute shaders, dispatched in batches across GPU cores.
    \item \textbf{State trie}: The Merkle Patricia Trie is stored in GPU-accessible memory with parallel hash computation for root updates.
\end{itemize}

\subsection{GPU Execution Pipeline}

\begin{figure}[t]
\centering
\begin{tikzpicture}[
    node distance=1.2cm,
    box/.style={rectangle, draw, minimum width=2.8cm, minimum height=0.8cm, align=center, fill=blue!10, font=\small},
    arrow/.style={->, thick}
]
    \node[box] (decode) {TX Decode\\(CPU)};
    \node[box, right=1.5cm of decode] (compile) {Opcode Compile\\(CPU $\to$ IR)};
    \node[box, right=1.5cm of compile] (dispatch) {GPU Dispatch\\(Batch)};
    \node[box, below=of dispatch] (execute) {Shader Execute\\(GPU Cores)};
    \node[box, left=1.5cm of execute] (state) {State Commit\\(Parallel Trie)};
    \node[box, left=1.5cm of state] (finalize) {Block Finalize\\(Quasar)};

    \draw[arrow] (decode) -- (compile);
    \draw[arrow] (compile) -- (dispatch);
    \draw[arrow] (dispatch) -- (execute);
    \draw[arrow] (execute) -- (state);
    \draw[arrow] (state) -- (finalize);
\end{tikzpicture}
\caption{Zoo EVM GPU execution pipeline. Transactions are batched for GPU dispatch after opcode compilation.}
\label{fig:pipeline}
\end{figure}

The pipeline processes transactions in three phases:

\begin{enumerate}
    \item \textbf{Preparation (CPU)}: Transaction decoding, signature verification, nonce checking, and opcode-to-IR compilation. This phase is parallelized across CPU cores.

    \item \textbf{Execution (GPU)}: IR programs are dispatched as compute shader workgroups. Each workgroup processes one transaction. Memory access conflicts are detected via atomic compare-and-swap on state slots. Conflicting transactions are re-executed sequentially.

    \item \textbf{Commitment (GPU+CPU)}: State trie updates are computed in parallel on GPU. The Merkle root is finalized on CPU and submitted to Quasar consensus.
\end{enumerate}

\subsection{Block Parameters}

\begin{table}[t]
\centering
\begin{tabular}{lcc}
\toprule
\textbf{Parameter} & \textbf{Zoo Network} & \textbf{Ethereum L1} \\
\midrule
Block gas limit & 1,000,000,000 & 30,000,000 \\
Target block time & 100ms & 12s \\
Max TXs/block (21K gas) & 47,619 & 1,428 \\
Consensus & Quasar + PoAI & Gasper (PoS) \\
Finality & 10ms (Quasar) & 12.8 min (2 epochs) \\
\bottomrule
\end{tabular}
\caption{Block parameters: Zoo Network vs Ethereum L1.}
\label{tab:params}
\end{table}

